\section*{Tables}

\begin{table}[h!]
\centering
\caption{Cross-species Bio-Gravitational Number $\mathcal{B}_g = EI/(MgL^2)$, computed directly from the tabulated $M$, $L$ and $EI$. $\mathcal{B}_g$ declines monotonically with body mass; on these stiffness estimates only the mouse exceeds $0.1$, and the human adult ($\mathcal{B}_g = 0.0101$) is \emph{not} distinguished from large quadrupeds by $\mathcal{B}_g$ alone. Shading marks the two human rows for reference, not a threshold crossing. $EI$ values are single-source estimates of heterogeneous provenance (see \texttt{Notes} in \texttt{data/species\_parameters.csv}); the ordering is provisional.}
\label{tab:cross_species}
\small
\begin{tabular}{lccccl}
\toprule
\textbf{Species} & \textbf{Mass (kg)} & \textbf{Length (m)} & \textbf{EI (Nm$^2$)} & \textbf{$\mathcal{B}_g$} & \textbf{Posture} \\
\midrule
Mouse & 0.03 & 0.05 & 0.00008 & 0.1087 & Quadruped \\
Rat & 0.3 & 0.15 & 0.002 & 0.0302 & Quadruped \\
Rabbit & 2.5 & 0.35 & 0.03 & 0.0100 & Quadruped \\
Cat & 4.0 & 0.45 & 0.06 & 0.0076 & Quadruped \\
Dog & 20.0 & 0.75 & 0.5 & 0.0045 & Quadruped \\
Chimpanzee & 50.0 & 0.5 & 1.6 & 0.0130 & Facultative biped \\
\rowcolor{yellow!15} Human (child) & 30.0 & 0.45 & 1.2 & 0.0201 & Biped \\
\rowcolor{yellow!15} Human (adult) & 70.0 & 0.6 & 2.5 & 0.0101 & Biped \\
Horse & 500.0 & 1.5 & 40.0 & 0.0036 & Quadruped \\
Elephant & 4000.0 & 2.5 & 550.0 & 0.0022 & Quadruped \\
\bottomrule
\end{tabular}
\end{table}

\begin{table}[h!]
\centering
\caption{Computational model parameters and biological justification}
\label{tab:parameters}
\small
\begin{tabular}{llll}
\toprule
\textbf{Parameter} & \textbf{Value} & \textbf{Units} & \textbf{Biological Basis} \\
\midrule
\multicolumn{4}{l}{\textit{Geometry}} \\
Length $L$ & 0.40 & m & Adult spine (sacrum to cranium) \\
Diameter $d$ & 0.03 & m & Typical vertebral body dimension \\
$n$ elements & 50--100 & -- & Convergence-tested discretization \\
\midrule
\multicolumn{4}{l}{\textit{Material Properties}} \\
Young's modulus $E_0$ & 1.0 & GPa & Effective modulus (bone + disc) \\
Area moment $I_{\mathrm{area}}$ & $1.0 \times 10^{-8}$ & m$^4$ & Cylindrical cross-section \\
Density $\rho$ & 1100 & kg/m$^3$ & Tissue-averaged density \\
Cross-section $A$ & $7.1 \times 10^{-4}$ & m$^2$ & $\pi (d/2)^2$ \\
\midrule
\multicolumn{4}{l}{\textit{IEC Coupling Parameters}} \\
$\chi_\kappa$ & 0--0.10 & m$^{-1}$ & Rest curvature coupling (sweep) \\
$\chi_E$ & 0.10 & -- & Stiffness modulation (fixed) \\
$\chi_M$ & 0--0.05 & N$\cdot$m & Active moment coupling \\
$\beta_1, \beta_2$ & 0.1, 0.05 & -- & Metric coupling constants \\
\midrule
\multicolumn{4}{l}{\textit{Information Field (Spinal Profile)}} \\
Cervical peak & $s/L = 0.80$ & -- & C1-C7 region \\
Lumbar peak & $s/L = 0.25$ & -- & L1-L5 region \\
Peak amplitude & 0.5--0.7 & -- & Normalized HOX expression \\
Baseline $I_0$ & 0.3 & -- & Background patterning \\
\midrule
\multicolumn{4}{l}{\textit{Loading \& Dynamics}} \\
Gravity $g$ & 0.01--1.0 & $g_\oplus$ & Microgravity to Earth \\
Damping $\nu$ & 0.1--1.0 & s$^{-1}$ & Numerical dissipation \\
Equilibrium criterion & $v_{\max} < 10^{-6}$ & m/s & Kinetic energy threshold \\
\bottomrule
\end{tabular}
\end{table}

% GENERATED by scripts/analysis/regenerate_table_thermodynamic.py — do not hand-edit.
% Source of truth: research/alphafold_v6_analysis/protein_metrics.json (AlphaFold DB v6 snapshot).
% Demand n=12 mean=3.08 (SD 1.50); Supply n=11 mean=1.79 (SD 0.59);
% MWU one-sided p=0.0105 two-sided=0.0210, Cohen's d=1.11, ratio=1.723.
% GENERATED by scripts/analysis/regenerate_table_thermodynamic.py — do not hand-edit.
% Source of truth: research/alphafold_v6_analysis/protein_metrics.json (AlphaFold DB v6 snapshot).
% Demand n=12 mean=3.08 (SD 1.50); Supply n=11 mean=1.79 (SD 0.59);
% MWU one-sided p=0.0105 two-sided=0.0210, Cohen's d=1.11, ratio=1.723.
\begin{table}[h!]
\centering
\caption{Thermodynamic cost proteins ($n=23$; 12 demand, 11 supply) classified as Demand (mechanosensory/cytoskeletal) or Supply (metabolic regulators), generated directly from the cached AlphaFold DB v6 metrics (\texttt{research/alphafold\_v6\_analysis/protein\_metrics.json}). Demand proteins are 72\% more elongated on average ($3.08$ vs $1.79$; Mann--Whitney $p = 0.021$ two-sided, $p = 0.011$ one-sided; Cohen's $d = 1.13$). Note the overlap: supply members PLOD1 and GHR individually exceed several demand members, and the gap survives FDR correction only marginally ($q = 0.049$, scored two-sided; Supplementary Table~S2). $^*$PIEZO2 v6 entry covers residues 1--709 of 2822; full-length anisotropy may differ.}
\label{tab:thermodynamic_cost_proteins}
\scriptsize
\begin{tabular}{llccllc}
\toprule
\textbf{Gene} & \textbf{UniProt} & \textbf{Category} & \textbf{Anisotropy} & \textbf{Disorder \%} & \textbf{pLDDT} & \textbf{Residues} \\
\midrule
\multicolumn{7}{l}{\textit{Demand-Side: Mechanosensory and Cytoskeletal (mean anisotropy $3.08 \pm 1.50$)}} \\
VIM & P08670 & Demand & 5.57 & 24.0 & 77.1 & 466 \\
PTK7 & Q13308 & Demand & 5.28 & 9.4 & 82.7 & 1070 \\
LMNA & P02545 & Demand & 4.71 & 25.8 & 76.4 & 664 \\
DSTYK & Q6XUX3 & Demand & 3.65 & 10.3 & 80.2 & 929 \\
PIEZO2$^{\star}$ & Q9H5I5 & Demand & 3.45 & 13.7 & 79.4 & 709 \\
PIEZO1 & Q92508 & Demand & 3.14 & 17.0 & 72.0 & 2521 \\
ADGRG6 & Q86SQ4 & Demand & 2.80 & 15.3 & 73.7 & 1221 \\
CAV1 & Q03135 & Demand & 2.52 & 2.8 & 78.4 & 178 \\
RUNX3 & Q13761 & Demand & 1.57 & 56.1 & 60.6 & 415 \\
EGR3 & Q06889 & Demand & 1.56 & 64.3 & 50.0 & 387 \\
FBN2 & P35556 & Demand & 1.42 & 12.3 & 68.4 & 1473 \\
LBX1 & P52951 & Demand & 1.36 & 33.3 & 61.7 & 348 \\
\midrule
\multicolumn{7}{l}{\textit{Supply-Side: Metabolic Regulators (mean anisotropy $1.79 \pm 0.59$)}} \\
PLOD1 & Q02809 & Supply & 3.13 & 3.0 & 92.7 & 727 \\
GHR & P10912 & Supply & 2.27 & 50.5 & 58.7 & 638 \\
ARNTL & O00327 & Supply & 2.13 & 40.4 & 65.5 & 626 \\
SHH & Q15465 & Supply & 1.94 & 16.5 & 78.4 & 462 \\
COMP & P49747 & Supply & 1.94 & 6.5 & 88.1 & 757 \\
CDKN1A & P38936 & Supply & 1.79 & 25.0 & 69.0 & 164 \\
IGF1R & P08069 & Supply & 1.44 & 15.7 & 78.0 & 1367 \\
PPARGC1A & Q9UBK2 & Supply & 1.35 & 61.9 & 52.7 & 798 \\
SOX9 & P48436 & Supply & 1.30 & 49.1 & 56.0 & 509 \\
SIRT1 & Q96EB6 & Supply & 1.28 & 46.6 & 65.0 & 747 \\
COL1A1 & P02452 & Supply & 1.12 & 67.3 & 52.7 & 1464 \\
\bottomrule
\end{tabular}
\end{table}

% GENERATED by scripts/analysis/regenerate_table_dissipation.py — do not hand-edit.
% Anisotropy/pLDDT from research/alphafold_v6_analysis/protein_metrics.json (v6 snapshot).
% MYLK/DMD/FLNA are not in the v6 snapshot (curated pre-v6 values, flagged in caption).
% GENERATED by scripts/analysis/regenerate_table_dissipation.py — do not hand-edit.
% Anisotropy/pLDDT from research/alphafold_v6_analysis/protein_metrics.json (v6 snapshot).
% MYLK/DMD/FLNA are not in the v6 snapshot (curated pre-v6 values, flagged in caption).
\begin{table}[h!]
\centering
\caption{The 16 candidate proteins analyzed in this study, grouped by mechanistic role ($\eta_p$ proprioceptive channels, $\eta_a$ active control, $\Gamma_m$ metabolic regulators) and characterized by AlphaFold-derived structural metrics (v6 snapshot; $^\dagger$ = curated value not present in the v6 export). The anisotropy ratio and morphology classification indicate the relative structural cost of maintaining directional sensing versus volumetric (scalar) function. Genes are prioritized by confidence-weighted scoring from the AFCC pipeline.}
\label{tab:dissipation_proteins}
\scriptsize
\begin{tabular}{llcclll}
\toprule
\textbf{Gene} & \textbf{UniProt} & \textbf{Category} & \textbf{Anisotropy} & \textbf{Morphology} & \textbf{pLDDT} & \textbf{Scaling} \\
\midrule
\multicolumn{7}{l}{\textit{Proprioceptive Channel Anchors ($\eta_{p}$)}} \\
PIEZO2 & Q9H5I5 & $\eta_{p}$ & 3.45 & Fibrous/Extended & 79.4 & $L$ \\
PIEZO1 & Q92508 & $\eta_{p}$ & 3.14 & Fibrous/Extended & 72.0 & $L^2$ \\
EGR3 & Q06889 & $\eta_{p}$ & 1.56 & Globular & 50.0 & $L$ \\
RUNX3 & Q13761 & $\eta_{p}$ & 1.57 & Globular & 60.6 & $L$ \\
\midrule
\multicolumn{7}{l}{\textit{Active Maintenance Anchors ($\eta_{a}$)}} \\
LMNA & P02545 & $\eta_{a}$ & 4.71 & Fibrous/Extended & 76.4 & $L^2$ \\
VIM & P08670 & $\eta_{a}$ & 5.57 & Fibrous/Extended & 77.1 & $L^3$ \\
LBX1 & P52951 & $\eta_{a}$ & 1.36 & Globular & 61.7 & $L^2$ \\
\midrule
\multicolumn{7}{l}{\textit{Metabolic Supply Scaling ($\Gamma_{m}$)}} \\
GHR & P10912 & $\Gamma_{m}$ & 2.27 & Intermediate & 58.7 & $L$ \\
ARNTL & O00327 & $\Gamma_{m}$ & 2.13 & Intermediate & 65.5 & $L$ \\
COL1A1 & P02452 & $\Gamma_{m}$ & 1.12 & Globular & 52.7 & $L^3$ \\
PPARGC1A & Q9UBK2 & $\Gamma_{m}$ & 1.35 & Globular & 52.7 & $L$ \\
SOX9 & P48436 & $\Gamma_{m}$ & 1.30 & Globular & 56.0 & $L$ \\
IGF1R & P08069 & $\Gamma_{m}$ & 1.44 & Globular & 78.0 & $L$ \\
\bottomrule
\end{tabular}
\end{table}
\begin{table}[h!]
\centering
\caption{Comprehensive Statistical Summary of the Information-Elasticity Coupling (IEC) Framework. All statistical tests were independently computed using \textit{scipy.stats} to summarize model-internal computational claims.}
\label{tab:statistical_summary}
\scriptsize
\begin{tabular}{p{3.5cm}p{3cm}p{4cm}p{3.5cm}}
\toprule
\textbf{Mechanism / Hypothesis} & \textbf{Metric Evaluated} & \textbf{Statistical Result} & \textbf{Significance Level} \\
\midrule
\multicolumn{4}{l}{\textit{Cross-Species Scaling \& The Allometric Trap}} \\
Allometric Scaling & $Bg$ vs. Body Mass ($n=10$) & Exponent = $-0.282$ (SE $0.058$) \newline $R^2 = 0.744$ \newline Bootstrap 95\% CI $[-0.347, -0.117]$ & $p = 1.31 \times 10^{-3}$ \\
\midrule
\multicolumn{4}{l}{\textit{Energy Deficit Window \& Scoliosis Onset}} \\
Critical Transition ($L_{crit}$) & Spine Length vs. simulated Cobb Angle & Monotone rise; $\mathrm{Cobb} \approx aL + bL^3$ & Deterministic sweep; no $p$-value is defined. \\
Structural anisotropy $\mathcal{A}$ & $L_{crit}$ vs. $\mathcal{A}$ & $r = -0.88$; $L_{crit}$ \emph{declines} with $\mathcal{A}$ \newline 20/50 points on search bounds & Deterministic sweep; no $p$-value is defined. \\
\midrule
\multicolumn{4}{l}{\textit{Mechanical Instability \& Directional-Stiffness Anisotropy}} \\
Directional Misalignment & Peak Stiffness Mismatch & Position = $0.596$ (Normalized) \newline Reduction = $31.1\%$ & Matching Thoracolumbar Apex (Verified) \\
\bottomrule
\end{tabular}
\end{table}
